\chapter{نتیجه‌گیری و کارهای آتی}
\newpage
\section{مروری بر روش پیشنهادی و تحلیل نتایج}



امروزه محاسبات کوانتومی به دلیل توانایی در افزایش سرعت الگوریتم‌ها توجه زیادی را به خود جلب نموده است. همچنین با پیشرفت روز افزون تکنولوژی، امروزه این محاسبات با محدودیت‌هایِ مانند اتلاف گرما و کاهش سرعت مواجه هستند. از طرفی قدرت محاسبات کوانتومی با افزایش تعداد کیوبیت‌های آن به صورت نمایی رشد می‌کند. یکی از چالش‌های اساسی در محاسبات کوانتومی این است که کیوبیت‌ها با کوچکترین اغتشاشی در محیط خراب شده و از بین می‌روند، مگر اینکه در دمای بسیار پایین نگه‌داری شوند. یکی از اغتشاشاتی که می توان منجر به از بین رفتن اطلاعات کوانتومی و ایجاد ناهمدوسی کیوبیت‌ها شود، تعامل آن‌ها در محیط است که با افزایش تعداد آن‌ها، این اطلاعات شکننده و در معرض خطا قرار می‌گیرند. بنابراین ساخت کامپیوترهای کوانتومی در مقیاس بزرگ و با تعداد کیوبیت‌های زیاد بسیار چالش برانگیز، و در عین‌حال دستیابی به آن، هدفی بسیار ارزشمند است.

محاسبات کوانتومی توزیع شده راه حل مناسبی برای ساخت سیستم‌های کوانتومی در مقیاس بزرگ می‌باشند. برای داشتن یک سیستم کوانتومی توزیع شده، نیاز به مکانیزمی برای ایجاد ارتباط بین زیر‌سیستم‌ها است. یکی از پروتکل‌های اصلی در محاسبات کوانتومی توزیع شده، پروتکل دورنوردی است که از خاصیت در‌هم‌تنیدگی کیوبیت‌ها برای توزیع نمودن اطلاعات کوانتومی استفاده می‌کند. این پروتکل دارای هزینه‌ی زیادی برای سیستم کوانتومی می‌باشد. بر طبق تئوری تکثیرناپذیری کوانتومی، زمانی که یک کیوبیت به مقصد خود دورنورد می‌شود،  نمی‌تواند در زیر سیستم خود استفاده گردد و تا زمان برگشت کیوبیت دورنوردی شده، امکان استفاده از کیوبیت در زیرسیستم خود وجود ندارد. به‌همین دلیل تعداد دورنوردی‌ها در یک سیستم کوانتومی توزیع شده چالش اساسی محسوب می‌گردد. در نتیجه ارائه یک الگوریتم جهت کاهش این تعداد، کمک قابل توجهی به طراحی و توسعه دیگر مدارهای کوانتومی خواهد نمود.

در این رساله، به طراحي سیستم کوانتومی توزیع شده با رويكرد کاهش دوبرری‌های کوانتومی پرداخته شد. از آنجایی که طراجی سیستم‌های کوانتومی با محدودیت‌های زیادی مواجه است، طراحی چنین سیستم با هدف کمینه‌سازی تعداد دورنوردی‌ها در یک فاز امری دشوار است. بنابراین طراحی و بهينه‌سازي  آن، در دو سطح انجام شد. در سطح اول به دنبال افرازبندی کیوبیت‌ها در تعدادی زیر سیستم کوانتومی با هدف توزیع متوزان آن‌ها صورت پذیرفت. در این سطح مدلی ریاضی برای مسئله افراز بندی کیوبیت‌ها ارائه گردید.  سپس جهت افراز بندي آن‌ها در دو زير سيستم و چند زير سيستم الگوريتم‌هاي جداگانه‌اي ارائه شد. سپس در سطح دوم با استفاده از افرازبندی و برچسب گذاری‌ بدست آمده جهت توزيع كيوبيت‌ها از سطح اول، سعی در کاهش تعداد دورنوردی‌های کوانتومی گردید. از آنجايي كه در سطح دوم زمان اجرايي دروازه‌ها امري مهم محسوب مي‌گردد، مدلي ماتريسي جهت نمايش دروازه‌هاي كوانتومي ارائه گرديد كه در آن ترتيب زمان اجراي دروازه‌ها و همچنين كنترل يا هدف بودن كيوبيت در آن نگهداري مي‌گردد. در این سطح، از این ایده الهام گرفته شد: زمانی‌که یکی از کیوبیت‌های یک دروازه‌ی سراسری از مبدا خود به مقصد مورد نظر دورنورد می‌گردد، می‌تواند قبل از برگشت آن به مبدا خود، توسط تعداد بیشتری دروازه‌های دیگر مورد استفاده قرار گیرد که این مسئله در سطح اول بررسی نمی‌شد. همین امر منجر به کاهش بیش از نیمی از تعداد دورنوردی‌های کوانتومی می‌گردد.


 
 \section{بررسي نتايج بدست آمده}
 
 روش پیشنهادی روی چندین مجموعه مدار کوانتومی از کتابخانه‌های مختلف تست شد. همچنین معیارهای متفاوتی جهت ارزیابی روش پیشنهادی، پیشنهاد گردید و با  الگوریتم‌های موجود مقایسه گردید. 
  
نتايجي كه از آزمايش‌هاي متعدد حاصل شد:
\begin{itemize}
\item
در سطح اول تعداد ارتباطات بدست آمده همان تعداد دروازه‌هاي سراسري بوده كه مي‌توان با اعمال بهينه‌سازي سطح دوم اين تعداد را به بيش از نصف كاهش داد.
\item
كاهش تعداد دورنوردي‌هاي بدست آمده در سطح دوم به اين معنا است كه مي‌توان تعداد زيادي از دروازه‌هاي سراسري را كه در سطح اول بدست آمده است، در سطح دوم به صورت محلي اجرا نمود. به عبارتي با دورنورد شدن يك كيوبيت از مبدا خود به مقصد مورد نظر، اين كيوبيت توسط تعداد زيادي دروازه‌ي سراسري ديگر با رعايت محدوديت‌هاي پيش نيازي مورد استفاده قرار گيرد.
\item
هر چه فاكتور توزان بار به صفر نزديك‌تر باشد، به اين معناست كه كيوبيت‌ها بايستي به طور متوازن‌تر در زير سيستم‌هاي كوانتومي توزيع شوند. در نتيجه تعداد دورنوردي‌هاي كوانتومي افزايش مي‌يابد. هرچه فاصله‌ی اين معيار از صفر بيشتر شود، الگوريتم توزيع، كيوبيت‌هاي كه بيشتر با يكديگر در تعامل هستند، در يك زيرسيتم قرار داده و این کیوبیت‌ها به طور محلی بیشتر با یکدیگر در ارتباط  بوده و در نتيجه تعداد دورنوردي‌هاي كوانتومي كاهش مي‌بابد.
 
\end{itemize}
\section{كارهاي آتي}
اگرچه نتايج حاصل از اين پژوهش رضايت بخش و قابل قبول مي‌باشد، با اين حال روزنه‌هاي جذاب و جالب توجهي پيش روي اين زمينه از پژوهش وجود دارد. از جمله اقداماتي كه در مسير افزايش كارايي روش پيشنهادي مي‌توان انجام داد، به اين موارد اشاره دارد:
\begin{itemize}
\item
اولين اقدام، طراحی یک سیستم توزیع شده با هدف کاهش تعداد دورنوردی‌های کوانتومی در یک سطح است به طوری که بتوان دو سطح را با یکدیگر ادغام نمود. 
\item
راهكار پيشنهادي ديگر، مي‌توان به اين نكته اشاره نمود كه به جاي جايگزين نمودن دروازه‌هاي با بيشتر از دو كيوبيت با دروازه‌هاي تك كيوبيتي و دو كيوبيتي، آن‌ها را جايگزين نکرده و عمل افرازبندي و بهينه‌سازي دورنوردي‌هاي كوانتومي روي مدارهایی با اين دروازه‌ها انجام گردد.
 \end{itemize}

